This chapter presents the architecture and methodology of the proposed L2M
system. The system is organized as a six-stage sequential pipeline in which
each stage produces a validated, typed output consumed by the next. Two of
the stages invoke an LLM; each of these is paired with a deterministic
fallback so that the pipeline always completes.

\section{Pipeline Overview}
\label{sec:mt-pipeline}

\ref{fig:pipeline} shows the six-stage pipeline. Raw lyrics are
parsed and normalized (Stage~1), analyzed for emotion and rhythm by the
first LLM call (Stage~2), and turned into a syllable-aligned note sequence
by the second LLM call (Stage~3). The result is converted into a typed
internal representation (Stage~4), exported to MIDI and MusicXML (Stage~5),
and optionally rendered to audio (Stage~6). At each LLM-dependent stage, a
fallback path activates automatically on API failure, timeout, or malformed
response.

\begin{figure}[ht]
\centering
\resizebox{0.78\textwidth}{!}{%
\begin{tikzpicture}[
  node distance = 0.45cm and 1.0cm,
  stage/.style={
    rectangle, rounded corners=4pt,
    draw=black!55, fill=gray!8, thick,
    minimum width=5.6cm, minimum height=0.72cm,
    align=center, font=\small
  },
  llmstage/.style={
    stage, fill=blue!10, draw=blue!55!black
  },
  optstage/.style={
    stage, draw=black!40, dashed, fill=gray!4
  },
  io/.style={
    ellipse, draw=black!60, thick, fill=white,
    minimum width=3.8cm, minimum height=0.58cm,
    align=center, font=\small
  },
  llmbadge/.style={
    rectangle, rounded corners=2pt,
    draw=blue!55!black, fill=blue!14,
    font=\scriptsize\sffamily, inner sep=2.5pt
  },
  fbbadge/.style={
    rectangle, rounded corners=2pt,
    draw=orange!65!black, fill=orange!14,
    font=\scriptsize\sffamily, inner sep=2.5pt
  },
  arr/.style={-{Stealth[length=5pt,width=4pt]}, thick, black!65},
  sidearr/.style={-{Stealth[length=3.5pt,width=3pt]}, thin, dashed},
  dtlabel/.style={font=\scriptsize\ttfamily, text=black!55, inner sep=1pt},
]

\node[io]                    (in)  {Lyrics Text};
\node[stage,   below=of in]  (s1)  {\textbf{Stage 1}\enspace Lyrics Parsing};
\node[llmstage,below=of s1]  (s2)  {\textbf{Stage 2}\enspace Emotion \& Rhythm Analysis};
\node[llmstage,below=of s2]  (s3)  {\textbf{Stage 3}\enspace Melody Structure Generation};
\node[stage,   below=of s3]  (s4)  {\textbf{Stage 4}\enspace Internal Representation};
\node[stage,   below=of s4]  (s5)  {\textbf{Stage 5}\enspace MIDI\,/\,MusicXML Export};
\node[optstage,below=of s5]  (s6)  {\textbf{Stage 6}\enspace Audio Rendering\;{\small\itshape(optional)}};

\draw[arr] (in) -- (s1);
\draw[arr] (s1) -- node[dtlabel,right=2pt]{\texttt{NormalizedLyrics}} (s2);
\draw[arr] (s2) -- node[dtlabel,right=2pt]{\texttt{EmotionAnalysis}}  (s3);
\draw[arr] (s3) -- node[dtlabel,right=2pt]{\texttt{MelodyStructure}}  (s4);
\draw[arr] (s4) -- node[dtlabel,right=2pt]{\texttt{Melody IR}}        (s5);
\draw[arr] (s5) -- node[dtlabel,right=2pt]{\texttt{.mid\;/\;.musicxml}} (s6);

\node[llmbadge, right=0.65cm of s2] (llm1) {LLM Call\,1};
\draw[sidearr, blue!55!black] (llm1.west) -- (s2.east);

\node[llmbadge, right=0.65cm of s3] (llm2) {LLM Call\,2};
\draw[sidearr, blue!55!black] (llm2.west) -- (s3.east);

\node[fbbadge, left=0.65cm of s2] (fb2) {Fallback};
\draw[sidearr, orange!65!black] (fb2.east) -- (s2.west);

\node[fbbadge, left=0.65cm of s3] (fb3) {Fallback};
\draw[sidearr, orange!65!black] (fb3.east) -- (s3.west);

\end{tikzpicture}
}% end resizebox
\caption[L2M six-stage pipeline]{L2M six-stage sequential pipeline. Blue
stages invoke LLM API calls; each has a deterministic fallback (orange).
Stage~6 (dashed) is optional and requires FluidSynth.}
\label{fig:pipeline}
\end{figure}

\section{Design Principles}
\label{sec:mt-design}

Five principles guided the system design; \ref{tab:design-principles}
summarizes their realization.

\begin{table}[ht]
\centering
\caption[Design principles and their realization]{Design principles and
their realization in L2M.}
\label{tab:design-principles}
\begin{tabular}{p{3.2cm} p{8.9cm}}
\toprule
\textbf{Principle} & \textbf{Realization} \\
\midrule
Modularity & Each stage is an independent service with defined input/output
types \\
Type safety & Pydantic~v2 models validate every inter-stage data structure \\
Robustness & Fallback heuristics at Stages 2 and 3 ensure a 100\% completion
rate \\
Observability & Structured logging with component-level context at every
stage \\
Extensibility & Service interfaces allow drop-in replacement (e.g., a
different LLM provider or audio synthesizer) \\
\bottomrule
\end{tabular}
\end{table}

Every transition in the data flow,
\begin{center}
\ttfamily\small
str $\rightarrow$ NormalizedLyrics $\rightarrow$ EmotionAnalysis
$\rightarrow$ MelodyStructure $\rightarrow$ Melody IR $\rightarrow$
music21 Stream $\rightarrow$ files,
\end{center}
is guarded by schema validation, catching malformed LLM outputs before they
propagate downstream.

\section{Stage 1: Lyrics Parsing}
\label{sec:mt-stage1}

The \texttt{LyricParser} service performs three preprocessing steps:

\begin{itemize}
    \item \textbf{Normalization}: collapse whitespace, strip excess
    punctuation, and handle line breaks;
    \item \textbf{Phrase segmentation}: split the text on newlines and
    sentence boundaries into lyric phrases;
    \item \textbf{Syllable estimation}: count vowel clusters per word, with
    adjustments for silent-\emph{e} endings, to obtain an approximate
    syllable count per phrase.
\end{itemize}

The heuristic syllable count serves two purposes: it provides a reference
against which the LLM's own count is cross-checked, and it drives the
fallback melody generator when the LLM path is unavailable.

\section{Stage 2: Emotion and Rhythm Analysis via LLM}
\label{sec:mt-stage2}

The first LLM call extracts the musical attributes of the lyrics using a
structured few-shot prompt. The response is a JSON object with four fields:
the dominant \texttt{emotion} (one of \emph{happy, sad, hopeful, tense,
calm, excited, melancholic}), a suggested \texttt{tempo} in BPM (validated
range 40--200), a \texttt{time\_signature} (e.g., \texttt{4/4},
\texttt{3/4}, \texttt{6/8}), and a per-line syllable breakdown in
\texttt{phrases}. The prompt is assembled from six components:

\begin{enumerate}
    \item System role assignment (``expert music composer and emotional
    analyst'');
    \item JSON output schema specification;
    \item Three diverse few-shot examples (happy, sad, hopeful);
    \item Empirically derived tempo range guidelines per emotion category;
    \item The target lyrics;
    \item A format enforcement instruction.
\end{enumerate}

An example response for the input \enquote{The sun will rise again}:

\begin{lstlisting}[style=jsonStyle]
{
  "emotion": "hopeful",
  "tempo": 90,
  "time_signature": "4/4",
  "phrases": [{"line": "The sun will rise again", "syllables": 6}]
}
\end{lstlisting}

\textbf{Fallback.} If the LLM call fails or returns invalid JSON, a
keyword-matching heuristic detects the emotion from a curated lexicon
(e.g., \emph{rain}, \emph{fade}, \emph{dark} $\rightarrow$ \emph{sad};
\emph{dance}, \emph{shine}, \emph{joy} $\rightarrow$ \emph{happy}) and
assigns the median tempo and the most common time signature for that emotion
category.

\section{Stage 3: Melody Structure Generation via LLM}
\label{sec:mt-stage3}

The second LLM call generates the syllable-aligned note sequence,
conditioned on the Stage~2 analysis. The prompt receives the detected
emotion, tempo, time signature, total syllable count $N$, and the lyrics,
and imposes the following constraints:

\begin{itemize}
    \item \textbf{Emotion-to-key guidance}: positive emotions
    $\rightarrow$ major keys; negative or tense emotions $\rightarrow$
    minor keys (per \ref{tab:emotion-mapping});
    \item \textbf{Hard alignment constraint}: \enquote{Generate EXACTLY $N$
    notes --- one per syllable};
    \item \textbf{Vocal range constraint}: all pitches within C3--C5;
    \item \textbf{Duration set}: $\{0.25, 0.5, 1.0, 2.0, 4.0\}$ beats;
    \item \textbf{Contour guidance}: avoid large melodic jumps; create
    natural phrase arcs.
\end{itemize}

The response is a JSON object containing the selected \texttt{key} and a
\texttt{melody} array of \texttt{\{note, duration, velocity\}} objects in
scientific pitch notation.

\subsection{Chunking for Long Lyrics}
\label{sec:mt-chunking}

When the total syllable count exceeds a configurable threshold (default:
30), the system splits the lyrics into phrase-level chunks and issues
sequential LLM calls. The last three notes of each chunk are carried into
the next prompt as context, maintaining local melodic continuity across
chunk boundaries.

\subsection{Emotion-to-Musical-Attribute Mapping}
\label{sec:mt-mapping}

\ref{tab:emotion-mapping} defines the explicit mapping used both to
guide the LLM (as prompt guidance) and to drive the fallback generator (as a
lookup table). The associations follow the music-psychology findings
reviewed in \Autoref{chap:background} \cite{hevner36expression,
juslin01musicemotion, gabrielsson10emotion}.

\begin{table}[ht]
\centering
\caption[Emotion-to-musical-attribute mapping]{Emotion-to-musical-attribute
mapping used for prompt guidance and fallback generation.}
\label{tab:emotion-mapping}
\begin{tabular}{l l c l}
\toprule
\textbf{Emotion} & \textbf{Key} & \textbf{Tempo (BPM)} & \textbf{Melodic
contour} \\
\midrule
Happy   & C major & 100--120 & Ascending \\
Hopeful & G major & 80--100  & Wavy (arch) \\
Sad     & A minor & 60--80   & Descending \\
Tense   & D minor & 90--110  & Erratic (high variance) \\
Calm    & F major & 60--80   & Balanced (small intervals) \\
Excited & D major & 120--140 & Ascending (steep) \\
\bottomrule
\end{tabular}
\end{table}

\subsection{Deterministic Fallback Generation}
\label{sec:mt-fallback}

If the melody-generation call fails, a deterministic generator constructs
the note sequence algorithmically from the emotion mapping.
Algorithm~\ref{alg:fallback} outlines the procedure. Four contour algorithms
are implemented:

\begin{itemize}
    \item \textbf{Ascending/Descending}: monotonic pitch progression with
    occasional plateaus based on syllable stress;
    \item \textbf{Wavy}: alternating up/down movement within a $\pm5$
    semitone range;
    \item \textbf{Balanced}: Gaussian-weighted note selection centered on
    the tonic;
    \item \textbf{Erratic}: uniform random jumps within scale degrees,
    conveying tension.
\end{itemize}

\begin{algorithm}[ht]
\caption{Deterministic fallback melody generation}
\label{alg:fallback}
\KwIn{Lyrics $L$ with $N$ estimated syllables; emotion $e$}
\KwOut{Melody $M$ with exactly $N$ notes}
$(key, [t_{min}, t_{max}], contour) \leftarrow \textsc{EmotionMap}(e)$\;
$tempo \leftarrow \operatorname{median}(t_{min}, t_{max})$\;
$scale \leftarrow \textsc{DiatonicScale}(key)$ restricted to C3--C5\;
$M \leftarrow \langle\rangle$\;
\For{$i \leftarrow 1$ \KwTo $N$}{
    $p_i \leftarrow \textsc{NextPitch}(scale, contour, M)$\;
    $d_i \leftarrow \textsc{DurationFor}(i, N)$
    \tcp*{phrase-final notes lengthened}
    append $(p_i, d_i)$ to $M$\;
}
\Return{$(key, tempo, M)$}
\end{algorithm}

Because the fallback is purely rule-based, it always terminates with a valid
melody, guaranteeing a 100\% pipeline completion rate regardless of network
or API conditions.

\section{Stage 4: Internal Representation}
\label{sec:mt-stage4}

The \texttt{MelodyGenerator} service converts the validated
\texttt{MelodyStructure} into a typed \texttt{Melody} dataclass --- the
internal representation (IR) --- comprising a list of \texttt{NoteEvent}
objects (scientific pitch, duration in quarter-note units, MIDI velocity)
together with the tempo, time signature, and key. Note-level validation
enforces pitch-range bounds and duration positivity before this stage
completes.

\section{Stage 5: MIDI and MusicXML Export}
\label{sec:mt-stage5}

The \texttt{MIDIWriter} service converts the \texttt{Melody} IR into a
music21 \texttt{Stream} object \cite{cuthbert10music21}, attaching a
metronome mark (tempo), key signature, time signature, and per-note
metadata (velocity, duration). The stream is exported to both MIDI
(\texttt{.mid}) and MusicXML (\texttt{.musicxml}) \cite{good01musicxml}
through music21's native exporters.

\section{Stage 6: Audio Rendering}
\label{sec:mt-stage6}

The optional \texttt{AudioRenderer} service invokes FluidSynth with a
configurable SoundFont (\texttt{.sf2}) to synthesize the MIDI file into PCM
audio (WAV, 44{,}100~Hz by default), with optional FFmpeg post-processing
for MP3 encoding.

\section{Implementation}
\label{sec:mt-implementation}

L2M is implemented as a Python package exposing both a command-line
interface and a programmatic API. \ref{tab:tech-stack} lists the
technology stack.

\begin{table}[ht]
\centering
\caption[Technology stack]{Technology stack of the L2M implementation.}
\label{tab:tech-stack}
\begin{tabular}{l l}
\toprule
\textbf{Component} & \textbf{Technology} \\
\midrule
Language & Python 3.9+ \\
LLM & OpenAI GPT-4o-mini (configurable) \\
Music notation & music21 \\
Data validation & Pydantic v2 \\
Audio synthesis & FluidSynth + FFmpeg \\
Interface & argparse CLI + Python API \\
Testing & pytest with coverage \\
\bottomrule
\end{tabular}
\end{table}

A single melody is generated from the command line as:

\begin{lstlisting}[style=jsonStyle, language=]
l2m --lyrics "The sun will rise again"
\end{lstlisting}

\noindent
The programmatic API mirrors the pipeline stages: a \texttt{LyricParser}
normalizes the input, an \texttt{LLMClient} performs the emotion analysis,
a \texttt{MelodyGenerator} produces the melody IR, and a \texttt{MIDIWriter}
exports the notation files. All configuration (model name, temperature,
maximum tokens, SoundFont path, sample rate) is supplied through environment
variables with sensible defaults.
